\subsection{TLB相关CSR寄存器的实现}

\begin{table}[h]
\centering
\caption{TLB相关CSR}
\begin{tabular}{|c|c|c|c|}
\hline
\textbf{CSR} & \textbf{主要更新源} & \textbf{冲突处理} & \textbf{优先级} \\
\hline
TLBIDX & CSR写、TLBSRCH、TLBRD & 无(只读部分) & TLBSRCH/TLBRD > CSR写 \\
\hline
TLBEHI & CSR写、异常、TLBRD & 重取指 & 异常 > CSR写 \\
\hline
TLBELO0/1 & CSR写、TLBRD & 无 & CSR写优先 \\
\hline
ASID & CSR写、TLBRD & 重取指 & TLBRD > CSR写 \\
\hline
CRMD(DA/PG) & CSR写、异常 & 重取指、特殊处理 & 异常 > CSR写 \\
\hline
TLBRENTRY & CSR写 & 无 & - \\
\hline
\end{tabular}
\end{table}

\subsubsection{TLBIDX寄存器的实现}

TLBIDX寄存器用于指定TLB操作的索引。其格式如下：

\begin{table}[h]
\centering
\caption{TLBIDX寄存器格式}
\begin{tabular}{|c|c|c|c|}
\hline
\textbf{位域} & \textbf{位数} & \textbf{含义} & \textbf{读写} \\
\hline
Index & 4 & TLB表项索引 & RW \\
\hline
PS & 6 & 页大小(只在TLBRD后更新) & R(via TLBRD) \\
\hline
NE & 1 & 无命中标志(TLBSRCH/TLBRD后更新) & R(via TLBSRCH/TLBRD) \\
\hline
\end{tabular}
\end{table}

在csr.v中，TLBIDX的实现包含多个更新来源：

\begin{lstlisting}[language=Verilog, caption=TLBIDX寄存器实现]
reg [31:0] csr_tlbidx;

always @(posedge clk) begin
    if (rst) begin
        csr_tlbidx <= 32'b0;
    end
    else if (csr_we && tlbidx_wsel) begin
        // CSR write: update Index field
        csr_tlbidx[`CSR_TLBIDX_INDEX] <= csr_wdata[`CSR_TLBIDX_INDEX];
    end
    else if (tlbsrch_en) begin
        // TLBSRCH: update Index and NE after search
        csr_tlbidx[`CSR_TLBIDX_INDEX] <= tlbsrch_index;
        csr_tlbidx[`CSR_TLBIDX_NE]    <= ~tlbsrch_found;
    end
    else if (tlbrd_en) begin
        // TLBRD: update PS and NE after read
        // If TLB entry is invalid (r_e=0), PS should be cleared
        csr_tlbidx[`CSR_TLBIDX_PS]    <= r_e ? r_ps : 6'b0;
        csr_tlbidx[`CSR_TLBIDX_NE]    <= ~r_e;
    end
end

// CSR read multiplexer
assign csr_rdata[`CSR_TLBIDX] = csr_tlbidx;
\end{lstlisting}

关键点：

\begin{enumerate}
    \item \textbf{NE位的清零保证}：当读取无效TLB项时（r\_e=0），PS域必须清为0，否则软件会读到错误的页大小信息。
    
    \item \textbf{多源更新优先级}：TLBSRCH和TLBRD操作的NE位更新优先级高于CSR写操作，确保查询/读取结果不被软件设置的值覆盖。
\end{enumerate}

\subsubsection{TLBEHI寄存器的实现}

TLBEHI寄存器存储虚拟页地址信息(VPPN)和其他页表项比较信息。

\begin{table}[h]
\centering
\caption{TLBEHI寄存器格式}
\begin{tabular}{|c|c|c|}
\hline
\textbf{位域} & \textbf{位数} & \textbf{含义} \\
\hline
VPPN & 19 & 虚拟页号(bit 31:13) \\
\hline
PS & 6 & 页大小(可选，某些实现用) \\
\hline
\end{tabular}
\end{table}

\begin{lstlisting}[language=Verilog, caption=TLBEHI寄存器实现]
reg [31:0] csr_tlbehi;

always @(posedge clk) begin
    if (rst) begin
        csr_tlbehi <= 32'b0;
    end
    else if (csr_we && tlbehi_wsel) begin
        // CSR write: software sets VPPN
        csr_tlbehi[`CSR_TLBEHI_VPPN] <= csr_wdata[`CSR_TLBEHI_VPPN];
    end
    else if (badv_ex) begin
        // Exception: hardware updates VPPN with bad virtual address
        csr_tlbehi[`CSR_TLBEHI_VPPN] <= badv_wdata[31:13];
    end
    else if (tlbrd_en) begin
        // TLBRD: read VPPN from TLB entry
        csr_tlbehi[`CSR_TLBEHI_VPPN] <= r_vppn;
    end
end

assign csr_rdata[`CSR_TLBEHI] = csr_tlbehi;
\end{lstlisting}

关键点：

\begin{enumerate}
    \item \textbf{异常时自动更新}：当发生TLB重填异常或其他地址相关异常时，硬件自动将坏虚地址的高19位写入VPPN。
    
    \item \textbf{异常优先于软件写}：异常触发的VPPN更新具有最高优先级，确保异常处理程序能获得正确的坏地址信息。
\end{enumerate}

\subsubsection{TLBELO0和TLBELO1寄存器的实现}

这两个寄存器分别存储TLB表项中的偶数页和奇数页的物理转换信息。

\begin{table}[h]
\centering
\caption{TLBELO寄存器格式}
\begin{tabular}{|c|c|c|}
\hline
\textbf{位域} & \textbf{位数} & \textbf{含义} \\
\hline
PPN & 20 & 物理页号(bit 31:12) \\
\hline
PLV & 2 & 特权等级(bit 11:10) \\
\hline
MAT & 2 & 存储访问类型(bit 9:8) \\
\hline
D & 1 & 脏位(bit 7) \\
\hline
V & 1 & 有效位(bit 6) \\
\hline
\end{tabular}
\end{table}

\begin{lstlisting}[language=Verilog, caption=TLBELO寄存器实现]
reg [31:0] csr_tlbelo0, csr_tlbelo1;

always @(posedge clk) begin
    if (rst) begin
        csr_tlbelo0 <= 32'b0;
        csr_tlbelo1 <= 32'b0;
    end
    else if (csr_we && tlbelo0_wsel) begin
        // CSR write: software sets page 0 translation info
        csr_tlbelo0[`CSR_TLBELO_PPN]  <= csr_wdata[`CSR_TLBELO_PPN];
        csr_tlbelo0[`CSR_TLBELO_PLV]  <= csr_wdata[`CSR_TLBELO_PLV];
        csr_tlbelo0[`CSR_TLBELO_MAT]  <= csr_wdata[`CSR_TLBELO_MAT];
        csr_tlbelo0[`CSR_TLBELO_D]    <= csr_wdata[`CSR_TLBELO_D];
        csr_tlbelo0[`CSR_TLBELO_V]    <= csr_wdata[`CSR_TLBELO_V];
    end
    else if (tlbrd_en) begin
        // TLBRD: read page 0 from TLB entry
        csr_tlbelo0[`CSR_TLBELO_PPN]  <= r_ppn0;
        csr_tlbelo0[`CSR_TLBELO_PLV]  <= r_plv0;
        csr_tlbelo0[`CSR_TLBELO_MAT]  <= r_mat0;
        csr_tlbelo0[`CSR_TLBELO_D]    <= r_d0;
        csr_tlbelo0[`CSR_TLBELO_V]    <= r_v0;
    end
    // TLBELO1 similar...
end

assign csr_rdata[`CSR_TLBELO0] = csr_tlbelo0;
assign csr_rdata[`CSR_TLBELO1] = csr_tlbelo1;
\end{lstlisting}

\subsubsection{ASID寄存器的实现}

ASID(Address Space Identifier)用于区分不同进程的虚地址空间，避免进程切换时清空整个TLB的性能损失。

\begin{table}[h]
\centering
\caption{ASID寄存器格式}
\begin{tabular}{|c|c|c|}
\hline
\textbf{位域} & \textbf{位数} & \textbf{含义} \\
\hline
ASID & 10 & 地址空间标识(bit 9:0) \\
\hline
\end{tabular}
\end{table}

\noindent\textbf{硬件实现}

\begin{lstlisting}[language=Verilog, caption=ASID寄存器实现]
reg [31:0] csr_asid;

always @(posedge clk) begin
    if (rst) begin
        csr_asid <= 32'b0;
    end
    else if (csr_we && asid_wsel) begin
        // CSR write: software updates ASID
        csr_asid[`CSR_ASID_ASID] <= csr_wdata[`CSR_ASID_ASID];
    end
    else if (tlbrd_en) begin
        // TLBRD: update ASID from TLB entry
        csr_asid[`CSR_ASID_ASID] <= r_asid;
    end
end

// ASID is read by TLB search ports directly
assign tlb_s1_asid = csr_asid[`CSR_ASID_ASID];
assign tlb_s0_asid = csr_asid[`CSR_ASID_ASID];

assign csr_rdata[`CSR_ASID] = csr_asid;
\end{lstlisting}

关键点：

\begin{enumerate}
    \item \textbf{直接连接TLB}：ASID寄存器的值直接驱动TLB查找模块的ASID输入，确保TLB能获得最新的ASID值进行查找。
    
    \item \textbf{重取指冲突}：由于取指和访存都会读取ASID用于TLB查找，而TLBRD可能修改ASID，此时采用重取指机制处理冲突。
\end{enumerate}

\subsubsection{TLBRENTRY寄存器的实现}

TLBRENTRY寄存器存储TLB重填异常处理的入口地址。与普通异常入口(EENTRY)不同，TLB重填有独立的入口。

\begin{lstlisting}[language=Verilog, caption=TLBRENTRY寄存器实现]
reg [31:0] csr_tlbrentry;

always @(posedge clk) begin
    if (rst) begin
        csr_tlbrentry <= 32'b0;
    end
    else if (csr_we && tlbrentry_wsel) begin
        // CSR write: software sets TLB refill exception entry
        csr_tlbrentry <= csr_wdata;
    end
end

// Used during TLB refill exception
assign tlbr_entry = csr_tlbrentry;

assign csr_rdata[`CSR_TLBRENTRY] = csr_tlbrentry;
\end{lstlisting}

\subsubsection{CRMD寄存器中的DA和PG位}

CRMD寄存器的DA(Direct translation)和PG(Paging)两位控制地址翻译模式：

\begin{itemize}
    \item DA=1, PG=0: 直接地址翻译模式(复位后的默认状态)
    \item DA=0, PG=1: 映射地址翻译模式(启用TLB和DMW)
    \item 其他组合: 保留或未定义
\end{itemize}

\noindent\textbf{硬件实现}

\begin{lstlisting}[language=Verilog, caption=CRMD的DA/PG位实现]
reg [31:0] csr_crmd;

always @(posedge clk) begin
    if (rst) begin
        // Reset to direct translation mode
        csr_crmd[`CSR_CRMD_DA]   <= 1'b1;
        csr_crmd[`CSR_CRMD_PG]   <= 1'b0;
    end
    else if (csr_we && crmd_wsel) begin
        // Software can write DA and PG
        csr_crmd[`CSR_CRMD_DA]   <= csr_wdata[`CSR_CRMD_DA];
        csr_crmd[`CSR_CRMD_PG]   <= csr_wdata[`CSR_CRMD_PG];
    end
    else if (wb_ecode == `ECODE_TLBR) begin
        // TLB refill exception: switch to direct translation
        csr_crmd[`CSR_CRMD_DA]   <= 1'b1;
        csr_crmd[`CSR_CRMD_PG]   <= 1'b0;
    end
    else if (csr_estat[`CSR_ESTAT_ECODE] == `ECODE_TLBR) begin
        // Exiting TLB refill: restore mapped translation
        csr_crmd[`CSR_CRMD_DA]   <= 1'b0;
        csr_crmd[`CSR_CRMD_PG]   <= 1'b1;
    end
end
\end{lstlisting}

关键点：

\begin{enumerate}
    \item \textbf{TLB重填异常处理}：当触发TLB重填异常时，硬件自动将DA置1、PG置0，进入直接地址翻译模式，避免异常处理程序自身再次触发TLB异常。
    
    \item \textbf{异常退出恢复}：执行ERTN指令返回TLB重填异常时，恢复到映射地址翻译模式。
    
    \item \textbf{重取指冲突}：由于DA/PG直接影响整个系统的地址翻译方式，CSR指令修改这两位会导致取指和访存的虚实转换规则改变。采用重取指机制处理。
\end{enumerate}

\subsubsection{CSR冲突的整体处理}

\noindent\textbf{CSR写后读冲突的种类}

在MMU中，CSR的写后读冲突主要包括：

\begin{enumerate}
    \item \textbf{显式写后读}：CSR写入指令(CSRWR/CSRXCHG)修改CSR，后续指令读取该CSR。
    
    \item \textbf{隐式写后读}：异常处理过程中硬件隐式修改CRMD、PRMD、ERA、ESTAT、BADV、TLBEHI等CSR，后续指令可能读取这些CSR。
\end{enumerate}

\noindent\textbf{冲突检测与处理}

在stage\_id.v中，对ID级指令进行CSR冲突检测：

\begin{lstlisting}[language=Verilog, caption=CSR写后读冲突检测]
// Detect explicit CSR write-after-read hazards
wire csr_hazard_ex  = ex_csr_we  & (ex_csr_num  == id_csr_num) & ex_valid;
wire csr_hazard_mem = mem_csr_we & (mem_csr_num == id_csr_num) & mem_valid;
wire csr_hazard_wb  = csr_we     & (csr_num     == id_csr_num) & wb_valid;

// Exception may implicitly write certain CSRs
wire ex_may_write_csr  = ex_ex_valid;
wire mem_may_write_csr = mem_ex_valid;
wire wb_may_write_csr  = wb_ex_valid;

// CSRs implicitly modified by exceptions
wire id_reads_ex_modified_csr = 
    (id_csr_num == `CSR_CRMD)   ||
    (id_csr_num == `CSR_PRMD)   ||
    (id_csr_num == `CSR_ERA)    ||
    (id_csr_num == `CSR_ESTAT)  ||
    (id_csr_num == `CSR_BADV)   ||
    (id_csr_num == `CSR_TLBEHI) ||
    (id_csr_num == `CSR_TLBIDX);

// Stall if there's a conflict
wire id_stall_csr = id_is_csr & 
    (csr_hazard_ex | csr_hazard_mem | csr_hazard_wb |
     (ex_may_write_csr & id_reads_ex_modified_csr) |
     (mem_may_write_csr & id_reads_ex_modified_csr) |
     (wb_may_write_csr & id_reads_ex_modified_csr));
\end{lstlisting}

\noindent\textbf{特殊冲突：重取指标志机制}

对于某些CSR修改（如CRMD、ASID）和TLB修改指令（TLBWR、TLBFILL、INVTLB），由于影响范围是整个地址翻译流水线，最早在 ID 级才能知道是TLB相关指令，此时 IF 级已经开始取址，无法简单地通过阻塞解决。采用重取指标志机制：

[TODO: fill it after implementation]

% \begin{lstlisting}[language=Verilog, caption=重取指标志设置]
% // Instructions that require refetch of subsequent instructions
% wire set_refetch_flag = wb_csr_we & (
%     (csr_num == `CSR_CRMD) |
%     (csr_num == `CSR_DMWO) |
%     (csr_num == `CSR_DMW1)
% ) | 
% wb_is_tlbwr | 
% wb_is_tlbfill | 
% wb_is_invtlb;
% 
% // Mark all younger instructions with refetch flag
% always @(posedge clk) begin
%     if (rst || flush) begin
%         refetch_flag <= 1'b0;
%     end
%     else if (set_refetch_flag) begin
%         // Set flag: subsequent instructions need refetch
%         refetch_flag <= 1'b1;
%     end
%     else if (refetch_instruction_reaches_wb) begin
%         // Clear flag after refetch instruction is processed
%         refetch_flag <= 1'b0;
%     end
% end
% \end{lstlisting}